\subsection{Mobile Application}

The mobile application is the primary user-facing interface for the Smart Water Filter Base. It is the client that displays live water-level readings for shared groups, handles device provisioning over Bluetooth Low Energy, manages group sharing, and surfaces low-water notifications on the user's phone. The web application is intentionally limited to account management and the setup guide; the mobile app is where day-to-day interaction with the system happens.

The mobile app is built with React Native using the Expo framework, and it shares its Supabase backend with the web application. The same email/password account works in both clients.

\subsubsection{System Architecture}
The mobile app interacts with the rest of the system in three ways:

\begin{itemize}
    \item \textbf{Supabase Auth} for sign-in, sign-up, and live profile state. The same Supabase project backs the web application.
    \item \textbf{Authenticated REST calls} to the Vercel-hosted API (e.g. \texttt{/api/groups}, \texttt{/api/devices}, \texttt{/api/profile}, \texttt{/api/app-state}, \texttt{/api/devices/ble-register}, \texttt{/api/groups/\{id\}/calibration}) using the Supabase access token as a Bearer credential. A default base URL is hard-coded in \texttt{src/api.js} and can be overridden with the \texttt{EXPO\_PUBLIC\_API\_BASE\_URL} environment variable.
    \item \textbf{Bluetooth Low Energy (BLE)} for the initial device-provisioning handshake with a Smart Base. The BLE path is currently gated to iOS; on Android, the Provision Device screen renders a notice explaining that QR + BLE provisioning is enabled for iOS first.
\end{itemize}

\noindent Low-water notifications are delivered as \emph{local} notifications through \texttt{expo-notifications}, scheduled by the mobile app itself rather than pushed from the server.

\subsubsection{Screens and User Flow}
The app is organized as a stack navigator (\texttt{@react-navigation/native-stack}) with the following screens, defined in \texttt{App.js}:

\begin{itemize}
    \item \textbf{Intro Screen.} Shown on first launch. Introduces the product before the user signs in. Completion is persisted in \texttt{AsyncStorage} under the \texttt{app:intro-completed} key so the intro is shown only once.
    \item \textbf{Auth Screen.} Email/password sign-in and sign-up, backed by Supabase Auth. Sign-up optionally captures a display name. The same credentials work in the web app.
    \item \textbf{Dashboard Screen.} The home screen after sign-in. It lists the user's groups and devices (fetched from \texttt{/api/groups} and \texttt{/api/devices}), lets the user create a new group, join an existing group by invite code, navigate into a group, open the device-provisioning flow, and sign out.
    \item \textbf{Provision Device Screen.} Walks the user through pairing a new base: scan the QR code on the base, enter Wi-Fi credentials, then transmit them to the base over BLE.
    \item \textbf{Group Screen.} Detailed view for a single group. Displays the latest reading for the group's associated device, exposes calibration controls, and is the destination opened when the user taps a low-water notification.
\end{itemize}

\subsubsection{Device Provisioning}
The Provision Device flow is the bridge between an out-of-the-box Smart Base and a cloud-connected device tied to the user's account. The flow is implemented in \texttt{ProvisionDeviceScreen.js} and proceeds in three stages:

\begin{enumerate}
    \item \textbf{Scan QR.} The user taps ``Scan ESP QR'', which opens the camera via \texttt{expo-camera} and decodes a QR sticker on the physical base. The screen accepts QR payloads as raw JSON, as URL query parameters (e.g. \texttt{espprov://...?device\_name=...}), or as base64-encoded JSON. Recognized fields include \texttt{device\_name}, \texttt{device\_id}, the BLE \texttt{service\_uuid} and characteristic UUIDs (\texttt{auth\_char\_uuid}, \texttt{wifi\_char\_uuid}, \texttt{status\_char\_uuid}), and metadata such as \texttt{pop}, \texttt{transport}, and \texttt{ver}. If UUID fields are omitted, fallback UUIDs hard-coded to match the \texttt{custom\_ble\_prov} firmware are used.
    \item \textbf{Register on the server.} A one-time \texttt{auth\_token} is generated on-device with a UUIDv4. When the user taps ``Send Token + Wi-Fi over BLE'', the app first calls \texttt{POST /api/devices/ble-register} with the \texttt{device\_id}, \texttt{auth\_token}, and chosen \texttt{device\_name}, so the backend will recognize the device when it later POSTs readings.
    \item \textbf{Transmit over BLE.} The app then uses \texttt{react-native-ble-plx} to scan for and connect to the advertising base, discover services and characteristics, request a 512-byte MTU, and write the payload in two parts: the \texttt{auth\_token} (base64-encoded UTF-8) to the auth characteristic, and the Wi-Fi credentials encoded as \texttt{ssid:password} (base64-encoded UTF-8) to the Wi-Fi characteristic. If a status characteristic UUID is provided, the app reads it back and displays the decoded value. After the writes, the app polls \texttt{/api/devices} for up to 30 seconds, waiting for the new \texttt{device\_id} to become visible from the server, then returns the user to the Dashboard.
\end{enumerate}

\noindent Because BLE is gated to iOS, the corresponding payloads are sent only when running on an iOS development build; on Android the button is disabled and the user sees an explanatory note.

\subsubsection{Groups and Sharing}
A user's bases are accessed through groups. The Dashboard supports creating a new group (optionally associating it with a device the user owns) and joining an existing group via an invite code. Group membership is read from \texttt{/api/groups}, and per-group details (including the latest reading) are read from \texttt{/api/groups/\{id\}}.

The Group screen also exposes calibration: empty and full reference points are captured by sending a \texttt{POST /api/groups/\{id\}/calibration} request with an \texttt{action} payload. Calibration is therefore stored per group (against the group's associated device) rather than as a single global record.

\subsubsection{Low-Water Notifications}
The mobile app monitors its own group devices for low-water conditions and surfaces local notifications when the level drops below threshold. Three pieces work together:

\begin{itemize}
    \item \textbf{\texttt{useLowWaterMonitor}.} A hook that runs while the user is signed in, the Supabase config is present, and the account is active. It polls \texttt{/api/groups} on an 8-second interval whenever the app is in the \texttt{active} state, then fetches \texttt{/api/groups/\{id\}} for each group that has a device attached and forwards the latest reading to the alerting helper.
    \item \textbf{\texttt{groupLowWaterAlerts}.} Computes the current water-level percentage and compares it against \texttt{LOW\_WATER\_THRESHOLD\_PERCENT = 20}. It tracks per-group alert state so that a single sub-threshold transition produces one notification, rather than firing repeatedly while the level stays low. State is cleared once the level returns at or above threshold (or on sign-out).
    \item \textbf{Expo Notifications.} When a fresh low-water transition is detected, the app schedules a local notification through \texttt{expo-notifications} carrying \texttt{\{ type: "low\_water", groupId, groupName \}} in its data payload. Tapping the notification is handled by \texttt{addNotificationResponseListener} and \texttt{handleInitialNotification} in \texttt{src/notifications.js}, which deep-link the user to the corresponding Group screen via the navigation reference held in \texttt{App.js}.
\end{itemize}

\noindent Because alerts are scheduled locally rather than pushed from the server, the system does not maintain a server-side device-token registry for the mobile app; alerts continue to work as long as the app has been in the foreground recently enough for the monitor to run.

\subsubsection{Account Deactivation}
\texttt{App.js} subscribes to Postgres change events on the signed-in user's row in the \texttt{profiles} table via Supabase Realtime. If \texttt{is\_active} flips to \texttt{false} (for example, from the web app), the mobile app immediately transitions into a deactivated state, clears all queued low-water notifications, and prompts the user to sign out. This keeps the two clients consistent: account-level controls live on the web, while live device interaction lives on mobile.
